% Generated draft: requires mathematical and source review.
\documentclass[11pt]{article}
\usepackage[margin=25mm]{geometry}
\usepackage{fontspec}
\setmainfont{Latin Modern Roman}
\usepackage{amsmath,amssymb,mathtools,unicode-math}
\setmathfont{Latin Modern Math}
\usepackage{enumitem,needspace,xurl,hyperref}
\setlength{\parindent}{0pt}
\setlength{\parskip}{5pt}
\setlength{\emergencystretch}{3em}
\newcommand{\N}{\mathbb{N}}
\newcommand{\Z}{\mathbb{Z}}
\newcommand{\Q}{\mathbb{Q}}
\newcommand{\R}{\mathbb{R}}
\newcommand{\C}{\mathbb{C}}
\newcommand{\F}{\mathbb{F}}
\newcommand{\A}{\mathbb{A}}
\newcommand{\PP}{\mathbb{P}}
\newcommand{\E}{\mathbb{E}}
\newcommand{\eps}{\varepsilon}
\newcommand{\dd}{\,\mathrm{d}}
\begin{document}
\section*{UnsolvedMath ID: 30002888}
\textbf{Source problem code:} catalogue entry
\subsection{Definitions and mathematical statement}
Let \(X,S,F\) denote the geometric space, stratification data, and coefficient data intended by the supplied source. Their admissible types and hypotheses are not specified in the record. Retain the source notation \(DS^{\mathrm{mix}}(X,F)\) and \(DS^b(X,F)\) for the mixed and bounded categories appearing in the question, and \(\operatorname{Perv}^{\mathrm{mix}}_S(X,F)\) and \(\operatorname{Perv}_S(X,F)\) for its mixed and ordinary perverse categories. These descriptions identify the intended roles of the symbols; the record does not supply their constructions.

\textbf{Source question.} For the data intended in the source, does there exist a forgetful functor
\[
\operatorname{For}\colon DS^{\mathrm{mix}}(X,F)\longrightarrow DS^b(X,F)
\]
that makes \(\operatorname{Perv}^{\mathrm{mix}}_S(X,F)\) a graded version of \(\operatorname{Perv}_S(X,F)\)?

\textbf{Unresolved qualifications.} The record does not define ``graded version,'' specify a grading shift or the required compatibility conditions for \(\operatorname{For}\), or state how the perverse categories relate to the displayed categories. It also does not specify the range of \((X,S,F)\) over which existence is sought. Consequently, this preserves the source question but does not determine a fully specified universally quantified existence statement. No particular geometric hypotheses, coefficient restrictions, categorical constructions, or grading axioms are imposed here.
\subsection{Short English statement}
The question asks whether a functor forgetting mixed structure can make mixed perverse sheaves into a graded version of ordinary perverse sheaves. The admissible geometric and coefficient data, the category constructions, and the meaning of ``graded version'' require clarification from the source. The supplied record labels the problem open; this draft does not independently verify that status.
\subsection{Sources}
\begin{itemize}
\item Supplied local UnsolvedMath registry record: permanent ID \(30002888\), source problem code \texttt{catalogue entry}, title \emph{Forgetful Functors for Mixed Perverse Sheaves}, in the Oberwolfach Reports Open Problems collection. No individual proposer is identified; the proposed-year field is \(2015\).
\item Supplied report citation: \emph{Enveloping Algebras and Geometric Representation Theory} (2016), Oberwolfach Reports, pp.\ 1385--1447, DOI \texttt{10.4171/owr/2015/25}. Source URL and supplied literature-source link, labeled ``Oberwolfach report 25'': \url{https://doi.org/10.4171/owr/2015/25}. The background describes these reports as documenting mathematical workshops and their open-problem discussions.
\item Literature assessment dated exactly \texttt{2026-08-21}: the supplied curation labels the problem open and states that the report explicitly describes existence of this forgetful functor as unknown in general, while later mixed-perverse results provide special cases without the stated general construction. This is supplied dated curation, not an independently verified assessment. The report above is the only listed literature source; no references for those later results are supplied. The background expressly requests independent verification of time-sensitive status and cited claims.
\item The registry labels the statement \texttt{corrected\_verified} and reports manual semantic, typographic, and LaTeX review confirming a meaning-preserving, self-contained cleanup without unsupported source-dependent additions. These are supplied editorial assertions; the unresolved definitions and hypotheses identified in this draft still require human mathematical review.
\end{itemize}
\end{document}
